\documentclass[11pt]{article}
\usepackage[margin=2.2cm]{geometry}
\usepackage{amsmath,amssymb}
\usepackage{booktabs}
\usepackage{hyperref}
\usepackage{parskip}

\title{Head-on test of BHLUMI + WarpX:\\ the beam-beam kick on a Bhabha lepton}
\author{BHLUMI\_WARPX repository, \texttt{benchmarks/} and \texttt{notebooks/01\_headon\_round\_benchmark.ipynb}}
\date{}

\begin{document}
\maketitle

\section*{1. The idea}

BHLUMI produces Bhabha events $e^+e^-\to e^+e^-(+n\gamma)$ for a perfect head-on collision. WarpX injects each
event at a vertex inside the colliding bunches and tracks the outgoing $e^\pm$ through the fields of the
\emph{opposite} bunch, which pull them slightly toward the beam axis. That kick is what biases the luminosity
measurement at FCC-ee.

To check the injection and tracking, we pick beams simple enough that the average kick can be computed by hand:
\begin{itemize}
  \item round Gaussian beams, $\sigma_x=\sigma_y=\sigma$, length $\sigma_z$, $N$ particles each, energy $E$ ($\gamma$);
  \item head-on (no crossing angle), no external field;
  \item rigid bunches: disruption $D = 2Nr_e\sigma_z/(\gamma\sigma^2)\ll 1$;
  \item no hourglass (zero emittance), no energy spread, no beamstrahlung (QED off).
\end{itemize}
The lepton is ultrarelativistic, leaves at a small polar angle $\theta$ and azimuth $\varphi$, and its own
bunch does not act on it (it moves almost parallel to the lepton, so its $E$ and $B$ forces cancel).

\section*{2. Derivation}

\textbf{Field.} The $e^+$ bunch moves at $-c$ with line density $\lambda(\zeta)=Ne\,g(\zeta)$,
$g(\zeta)=e^{-\zeta^2/2\sigma_z^2}/(\sqrt{2\pi}\,\sigma_z)$. Its field is transverse:
\begin{equation}
\mathbf E_\perp=\frac{\lambda(\zeta)}{2\pi\varepsilon_0}\,\frac{\mathbf r}{r^2}\left(1-e^{-r^2/2\sigma^2}\right),
\qquad c\mathbf B_\perp=\hat z\times\mathbf E_\perp .
\end{equation}

\textbf{Force.} For a lepton moving at $+c$, the electric and magnetic forces add: $\mathbf F_\perp=2q\mathbf E_\perp$.

\textbf{Path.} Lepton and bunch approach at $2c$. After the lepton has crossed a length $u=\zeta-\zeta_0$ of the
bunch, $t=u/2c$ and its transverse position is $\mathbf r(u)=\mathbf r_0+\hat n\,\theta u/2$, with
$\hat n=(\cos\varphi,\sin\varphi)$.

\textbf{Kick for one event.} Integrating the force along the path, projecting on $\hat n$ and dividing by
$p=\gamma' m c$ (with $r_e=e^2/4\pi\varepsilon_0 mc^2$):
\begin{equation}
\Delta\theta=-\frac{2Nr_e}{\gamma'}\int_0^\infty g(\zeta_0+u)\,
\frac{\mathbf r\cdot\hat n}{r^2}\left(1-e^{-r^2/2\sigma^2}\right)du .
\end{equation}

\textbf{Average over the transverse vertex.} Vertices follow the product of the two beam densities, so
$\mathbf r_0$ is Gaussian with variance $\sigma^2/2$ per axis. Averaging the field point over that offset gives the
field of a Gaussian with $\sigma_{\rm eff}^2=\sigma^2+\sigma^2/2=\tfrac32\sigma^2$:
\begin{equation}
\left\langle\frac{\mathbf r\cdot\hat n}{r^2}\left(1-e^{-r^2/2\sigma^2}\right)\right\rangle
=\frac{2}{\theta u}\left(1-e^{-\theta^2u^2/12\sigma^2}\right).
\end{equation}

\textbf{Average over the longitudinal vertex.} $\zeta_0\sim g$, and
$\int g(\zeta_0)\,g(\zeta_0+u)\,d\zeta_0=e^{-u^2/4\sigma_z^2}/(2\sqrt\pi\,\sigma_z)$. Then
\begin{equation}
\langle\Delta\theta\rangle=-\frac{2Nr_e}{\sqrt\pi\,\gamma'\theta\,\sigma_z}
\int_0^\infty\frac{e^{-bu^2}-e^{-(a+b)u^2}}{u}\,du,
\qquad a=\frac{\theta^2}{12\sigma^2},\quad b=\frac{1}{4\sigma_z^2},
\end{equation}
and this (Frullani) integral equals $\tfrac12\ln\!\big((a+b)/b\big)$.

\section*{3. Result}

\begin{equation}
\boxed{\;\langle\Delta\theta\rangle(\theta)=-\frac{N r_e}{\sqrt{\pi}\,\gamma'\,\theta\,\sigma_z}
\ln\!\left(1+\frac{\theta^2\sigma_z^2}{3\sigma^2}\right)\;}
\end{equation}
\begin{itemize}
  \item The minus sign: the lepton is pulled toward the axis. The same formula holds for the $e^+$ ($\theta$ from $-z$).
  \item It does not depend on $\varphi$ for round beams.
  \item It scales as $1/\gamma'$, the energy of the \emph{outgoing} lepton, so it is evaluated per event with
        BHLUMI's lepton energy.
  \item Limits: $\theta\sigma_z\ll\sigma$ gives a thin lens, $\propto\theta$;
        $\theta\sigma_z\gg\sigma$ gives a line charge, $\propto\ln\theta/\theta$.
\end{itemize}

\textbf{Default input} (\texttt{benchmarks/headon\_round/inputs\_headon\_round.txt}): $E=45.6$\,GeV,
$N=2\times10^{10}$, $\sigma=10\,\mu$m, $\sigma_z=1$\,mm, so $D=0.013$ and
$\mathcal L_{\rm BX}=N^2/4\pi\sigma^2=3.18\times10^{29}\,\mathrm{m^{-2}}$.

\begin{center}
\begin{tabular}{cccccc}
\toprule
$\theta$ [mrad] & 40 & 60 & 80 & 100 & 120 \\
$\langle\Delta\theta\rangle$ [$\mu$rad] & $-16.44$ & $-15.23$ & $-13.84$ & $-12.60$ & $-11.56$ \\
\bottomrule
\end{tabular}
\end{center}
A direct numerical average over vertices and trajectories agrees within its statistical error
(\texttt{bhabha\_theory.kick\_mc}); the per-event spread is about 7--9\,$\mu$rad rms.

\section*{4. What the test compares}

For each injected event, WarpX gives the lepton's final angle; the test compares
$\Delta\theta=\theta_{\rm out}-\theta_0$, averaged in bins of $\theta_0$, with the formula above evaluated for
the same events. The theory reads every number from the WarpX input file, so the two cannot drift apart. Before
comparing, the notebook checks that the file meets the assumptions: round, head-on, $D<0.05$, no emittance, no QED,
no external field, at least 5 cells per $\sigma$, at least 10 steps over the cutoff length $2\sqrt3\,\sigma/\theta$,
a box wide enough that truncating the kick costs $<10^{-3}$, and a run long enough for the last vertex.

Only after this test agrees do we add, one at a time, beamstrahlung, flat beams and the FCC-ee crossing angle.

\end{document}
